% Part: normal-modal-logic
% Chapter: filtrations
% Section: S5-fmp

\documentclass[../../../include/open-logic-section]{subfiles}

\begin{document}

\olfileid{nml}{fil}{fmp}

\olsection{\Log{K} 与 \Log{S5} 具有有限模型性}

\begin{defn}
  一个模态逻辑系统 $\Sigma$ 具有\emph{有限模型性}，如果每当公式~$!A$ 在 $\Sigma$ 的某个模型的某个世界上为真时，$!A$ 在 $\Sigma$ 的某个\emph{有限}模型的某个世界上也为真。
\end{defn}

\begin{prop}\ollabel{prop:K-fmp}
  \Log{K} 具有有限模型性。
\end{prop}

\begin{proof}
  \Log{K} 是全体有效公式的集合，亦即任何模型都是~\Log{K} 的模型。
  根据 \olref[fil]{thm:filtrations}，如果
  $\mSat{M}{!A}[w]$，那么 $\mSat{M^*}{!A}[w]$ 对 $\mModel{M}$ 经由 $!A$ 的子公式集 $\Gamma$ 所作的任何过滤成立。
  任何公式都只有有限多个子公式，故 $\Gamma$ 是有限的。
  由 \olref[fin]{prop:filt-are-finite} 可知，$\card{W^*} \le 2^n$，其中 $n$ 是~$\Gamma$ 中公式的个数。
  又因为 \Log{K} 对模型没有限制，所以 $\mModel{M^*}$ 是一个 \Log{K}-模型。
\end{proof}

为了通过过滤证明一个逻辑~\Log{L} 具有有限模型性，本质在于确保一个 \Log{L}-模型的过滤本身也是 \Log{L}-模型。
这通常需要相当多的工作，且并非任何过滤都能产生 \Log{L}-模型。然而，对于全通模型，这一点仍然成立。

\begin{prop}\ollabel{prop:univ-fin}
  令 $\mClass{U}$ 为全通模型的类（参见 \olref[frd][es5]{prop:S5=univ}），$\mClass{U}_\mathrm{Fin}$ 为所有有限全通模型的类。
  那么，公式~$!A$ 在 $\mClass{U}$ 中有效，当且仅当它在 $\mClass{U}_\mathrm{Fin}$ 中有效。
\end{prop}

\begin{proof}
  有限全通模型是全通模型，因此从左到右的方向是平凡的。
  对于从右到左的方向，假设~$!A$ 在某个全通模型 $\mModel{M}$ 的某个世界 $w$ 处为假。
  令 $\Gamma$ 包含 $!A$ 及其所有子公式；显然 $\Gamma$ 是有限的。
  取 $\mModel{M}$ 的一个过滤 $\mModel{M^*}$；由 \olref[fin]{prop:filt-are-finite}，$\mModel{M^*}$ 是有限的，
  且由 \olref[fil]{thm:filtrations}，$!A$ 在 $\mModel{M^*}$ 中的 $[w]$ 处为假。
  只需注意到 $\mModel{M^*}$ 也是全通的：给定 $u$ 和 $v$，根据假设有 $Ruv$，再由 \olref[fil]{defn:filtration}\olref[fil]{defn:filtration-R}，亦有 $R^*[u][v]$。
\end{proof}

\begin{cor}\ollabel{cor:S5fmp}
  \Log{S5} 具有有限模型性。
\end{cor}

\begin{proof}
  由 \olref[frd][es5]{prop:S5=univ}，如果 $!A$ 在某个自反且欧几里得的模型的某个世界上为真，那么它在某个全通模型的某个世界上也为真。
  由 \olref{prop:univ-fin}，它在某个有限全通模型（即对该模型经由 $!A$ 的子公式集进行过滤所得的模型）的某个世界上为真。
  每个全通模型也都自反且欧几里得；因此 $!A$ 在某个有限的自反且欧几里得的模型的某个世界上为真。
\end{proof}

\begin{prob}
  证明：持续（或自反）模型的任何过滤也分别是持续（或自反）的。
\end{prob}

\begin{prob}
  找出一个对称（或传递，或欧几里得）模型的非对称（或非传递，或非欧几里得）过滤。
\end{prob}

\end{document}